% =============================================================================
% STUDY DOSSIER — MOLAMBUDOS: THE DIARY OF PATIENT 1,260
% Critical apparatus. NOT part of the fiction; not to be bound with the book.
% =============================================================================
\documentclass[11pt,a4paper]{article}

\usepackage[T1]{fontenc}
\usepackage[utf8]{inputenc}
\usepackage[english]{babel}
\usepackage{ebgaramond}
\usepackage[margin=2.6cm]{geometry}
\usepackage{xcolor}
\usepackage{booktabs}
\usepackage{enumitem}
\usepackage{titlesec}
\usepackage[hidelinks]{hyperref}

\definecolor{sangue}{HTML}{8B0000}
\definecolor{tinta}{HTML}{2E241A}
\titleformat{\section}{\Large\bfseries\color{sangue}}{\thesection}{0.7em}{}
\titleformat{\subsection}{\large\bfseries\color{tinta}}{\thesubsection}{0.7em}{}
\setlist{itemsep=0.25em, topsep=0.35em}

\newcommand{\caveat}[1]{%
  \par\vspace{0.5em}\noindent\fcolorbox{sangue!40}{sangue!4}{%
  \parbox{0.965\linewidth}{\small\textbf{Methodological caveat.} #1}}\par\vspace{0.5em}}

\title{\vspace{-1.5cm}\textbf{Molambudos --- The Diary of Patient 1,260}\\[0.3em]
\large Study dossier: literary project, techniques, theory, limits}
\author{\small Critical apparatus to the trilingual edition}
\date{\small 5 August 2026}

\begin{document}
\maketitle
\thispagestyle{empty}

\begin{abstract}
\noindent This dossier describes and analyses \textit{Molambudos --- The Diary
of Patient 1,260}, an archive-novel built as printed hypertext and issued in
three language editions plus a trilingual one. It gathers the work's measured
technical record, its formal architecture, the literary lineage it belongs to,
the narrative techniques it deploys, the theoretical frameworks it mobilises,
its historical basis, and a critical analysis that includes verified
limitations. This is a study apparatus, not promotional copy: sections
\ref{sec:limits} and \ref{sec:validation} exist to state what the work has
\textit{not} achieved and what has \textit{not} been externally validated.
\end{abstract}

\caveat{Every figure in this dossier was measured against the work's source
files on 2026-08-05, not estimated. Where a claim is interpretive, the text
says so. Where a claim would require external validation that does not exist
--- reception, impact, comparative merit --- the dossier abstains rather than
guesses. See \S\ref{sec:validation}.}

% =============================================================================
\section{Technical record}

\begin{center}\small
\begin{tabular}{@{}ll@{}}
\toprule
\textbf{Title} & Molambudos --- The Diary of Patient 1,260 \\
\textbf{Author} & Marcelo Dias de Carvalho Filho \\
\textbf{Co-author and illustrator} & Marcelo Claro \\
\textbf{ISBN} & 9798189170492 \\
\textbf{Genre} & Archive-novel; documentary psychological horror \\
\textbf{Form} & Printed hypertext, 84 fragments, 216 reading routes \\
\textbf{Languages} & Portuguese (BR), English, Simplified Chinese \\
\textbf{Print format} & 160 $\times$ 230 mm, paperback \\
\bottomrule
\end{tabular}
\end{center}

\vspace{0.6em}
\noindent\textbf{Measured extent by edition}

\begin{center}\small
\begin{tabular}{@{}lrrr@{}}
\toprule
\textbf{Edition} & \textbf{Digital} & \textbf{Print 160$\times$230} & \textbf{Words} \\
\midrule
Portuguese & 415 pp. & 435 pp. & 43,730 \\
English    & 411 pp. & 427 pp. & 44,978 \\
Chinese    & 387 pp. & 397 pp. & 70,750 characters \\
Trilingual & 1,061 pp. & 1,115 pp. & --- \\
\bottomrule
\end{tabular}
\end{center}

\noindent\small Word counts are not comparable between Chinese and Western
languages; for Chinese the character count is given, which is the usual metric.

% =============================================================================
\section{The literary project}

\textit{Molambudos} presents itself as a \textbf{forensic archive}, not as a
novel. The reader receives 84 fragments --- memories, medical reports,
letters, audio transcripts, clinical forms --- assembled by an anonymous
archivist around a single object: the diary of Joaquim Antônio Correia,
registered at the Barbacena Colony Hospital under the number 1,260.

The declared project, legible in both structure and paratext, has three
movements:

\begin{enumerate}
\item \textbf{Recovering a buried history.} The Barbacena Colony existed;
across the twentieth century, historiography estimates that tens of thousands
of people died within its walls, many with no psychiatric diagnosis at all.
The work is built on that ground.

\item \textbf{Making form carry the argument.} Fragmentation is not ornament:
it enacts the impossibility of reconstructing a life from documentary
remains. The reader must assemble the case, and the assembly is always partial.

\item \textbf{Implicating the reader.} The work refuses the spectator
position. An entire family of fragments (CONT) addresses the reader in the
second person and inscribes them in the chain of transmission as
\textbf{Patient 1,263}.
\end{enumerate}

\subsection{The chain of transmission}

The central device is a numbered queue of occupied bodies:

\begin{center}\small
\begin{tabular}{@{}cll@{}}
\toprule
\textbf{No.} & \textbf{Who} & \textbf{Where fixed} \\
\midrule
1,259 & the one before Joaquim & CONT-07, CONT-08 \\
1,260 & Joaquim Antônio Correia & title of the work \\
1,261 & Dr Heitor Oliveira & DOC-16 (signs it) \\
1,262 & Dr Lúcia Mendes & LUC-10, LUC-11, Epilogue \\
\textbf{1,263} & \textbf{the reader} & LUC-12, DOC-24, CONT-09, protocol \\
\bottomrule
\end{tabular}
\end{center}

The precision of this count is a condition of the device working: it is what
turns the last position into a personal threat rather than a vague metaphor.

% =============================================================================
\section{Formal architecture}

\subsection{The four fragment families}

\begin{center}\small
\begin{tabular}{@{}llrrr@{}}
\toprule
\textbf{Code} & \textbf{Nature} & \textbf{N} & \textbf{Words} & \textbf{Mean} \\
\midrule
MEM  & Memories of Patient 1,260 & 27 & 13,310 & 492 \\
DOC  & Documents, reports, letters & 27 & 15,247 & 564 \\
LUC  & Lúcia Mendes's investigation (2026) & 17 & 8,934 & 525 \\
CONT & Reader contamination protocols & 13 & 6,239 & 479 \\
\midrule
     & \textbf{Total} & \textbf{84} & \textbf{43,730} & 521 \\
\bottomrule
\end{tabular}
\end{center}

The variation in length is wide and functional: from 130 words (MEM-26) to
2,310 (DOC-03, the autopsy report). Short fragments land as blows; long ones
work as documentary immersion.

\subsection{Five parts, three acts}

Linear reading runs through five parts --- Sertão (9), Colony (16), Oliveira's
Diary and Reports (25), Lúcia's Investigation (20), Contamination (14 +
Epilogue) --- grouped into three acts: Trauma (9 fragments),
Institutionalisation (41) and The Cycle (35).

The asymmetry is meaningful: the central act, corresponding to 62 years of
confinement, occupies nearly half the work. The duration of the confinement is
reproduced in the duration of the reading.

\subsection{The route network}

\begin{center}\small
\begin{tabular}{@{}lr@{}}
\toprule
Fragments & 84 \\
Reading routes & 216 \\
Distinct destinations & 85 (84 fragments + Epilogue) \\
Fragments with no inbound route & 0 \\
Fragments from which the Epilogue is reachable & 84 / 84 \\
Cycles (strongly connected components $>$ 1) & 6 \\
Cycles with no exit & 0 \\
Most referenced node & DOC-01, with 20 inbound routes \\
\bottomrule
\end{tabular}
\end{center}

Two properties are deliberate and were verified: \textbf{every route leads to
the Epilogue} and \textbf{no cycle traps the reader}. The cycles are preserved
--- they are the work's thesis, which states that ``the cycle begins again''
--- but none is a trap: from any fragment there is a finite path to the end.

\caveat{The alternative was a strictly acyclic graph. It was measured and
rejected: it would cost 85 of the 216 routes, precisely the
document$\to$memory cross-references that give the work its archival
character, and it would contradict the text's own thesis. The choice is the
author's and is on record.}

% =============================================================================
\section{Literary and narrative lineage}

\caveat{This section identifies \textbf{formal affinities verifiable in the
text}. It does not claim the author read the works cited or consciously
derived from them. Lineage here means kinship of procedure --- what criticism
can compare --- not a genealogy of reading.}

\subsection{The archive-novel and the forged document}

The work belongs to the tradition of literary found footage, whose modern
matrix is epistolary: \textit{Dracula} (1897) assembles horror from diaries,
letters and clippings, and delegates to the reader the synthesis no character
achieves. \textit{Molambudos} adopts the same principle and radicalises it by
including forensic reports, chains of custody and conservation assessments ---
bureaucratic genres, not intimate ones.

It also stands near \textit{House of Leaves} (Danielewski, 2000) in its use of
the page as space, its navigation apparatus and its bibliographic vertigo. It
differs on one decisive point: where Danielewski multiplies levels of
commentary into a labyrinth, \textit{Molambudos} keeps the chain short and
named, and uses it to reach the reader.

\subsection{Hypertext fiction on a printed support}

\textit{Hopscotch} (Cortázar, 1963) established the precedent of a book with
more than one reading order. \textit{Molambudos} goes beyond Cortázar's
``board'' by offering a dense network (216 edges over 85 nodes) with
properties guaranteed by automated verification --- convergence and absence of
entrapment.

\caveat{That verification is a feature of \textit{production}, not of literary
merit. It guarantees the navigational promise is kept; it says nothing about
the quality of the prose.}

\subsection{Testimony and atrocity literature}

The work converses with the Brazilian tradition of denouncing the asylum
system, whose contemporary journalistic landmark is \textit{Holocausto
Brasileiro} (Daniela Arbex, 2013), cited by name in DOC-19. In literary
register, its first part echoes \textit{Barren Lives} (Graciliano Ramos,
1938): the drought, the family's march, the stripped-down language.

The underlying theoretical question is that of \textit{representability}: how
to narrate an atrocity without spectacularising it. The work's choice ---
supernatural horror as an explicit metaphor, declared as such in the
\textit{Historical Note} --- is a position, and a debatable one, examined in
\S\ref{sec:critique}.

\subsection{Second person and the implicated reader}

Second-person address has a canonical precedent in \textit{If on a winter's
night a traveller} (Calvino, 1979), where it is playful and metaliterary.
\textit{Molambudos} inverts the key: not to play with reading, but to accuse
it. Twenty-seven of the 84 fragments sustain the ``you''.

% =============================================================================
\section{Narrative techniques}

\subsection{Induction and suggestion}

The CONT fragments deploy identifiable procedures from the literature of
clinical suggestion:

\begin{itemize}
\item \textbf{Double bind.} CONT-03 builds an alternative with no exit:
``\textit{If yes: the smell is there. [\ldots] If no: you are lying. You
lifted it. Everyone lifts it.}'' Both answers confirm the premise.

\item \textbf{Ideomotor suggestion.} ``\textit{Later, when you are distracted,
your finger will return to the page on its own.}''

\item \textbf{Post-hypnotic suggestion.} ``\textit{It will be on your pillow
tonight.}''

\item \textbf{Physiological anchoring.} Describing symptoms (breathing,
pupils, extremity temperature) tends to induce the attention that confirms them.

\item \textbf{Directed pareidolia.} ``\textit{Look at the upper right corner of
the room [\ldots] Fix your eye on the dark corner for 10 seconds.}''
\end{itemize}

\subsection{The inversion of absence}

The work's most sophisticated technique appears in CONT-04, ``Night 4'': after
three nights of nightmare, nothing happens. And the text converts the relief
into threat: ``\textit{the absence of fear was worse than the fear [\ldots]
The waiting is the true contagion.}''

The procedure is notable because it removes the stimulus and keeps the effect
--- transferring the production of horror from the text to the reader.

\subsection{The clinical instrument as form}

DOC-09 and CONT-05 take the shape of psychiatric self-assessment scales, with
items, options and a final score. Item 15 of DOC-09 --- ``\textit{If you could
stop reading now and never open this book again, would you stop? ( ) Yes ( )
No ( ) I already tried}'' --- works by double meaning: within the fiction it
refers to reading; in the register of the clinical instrument it imitates, it
reads as a suicidality item.

\caveat{That ambiguity is effective, and it is also why the front warning page,
with its crisis line, is not ornament. See \S\ref{sec:ethics}.}

\subsection{Marginalia and diegetic paratext}

The work uses marginal notes that comment on the text from outside, patient
records, archive seals and a fully diegetic ``Caution'' page. The paratext
participates in the fiction rather than framing it.

% =============================================================================
\section{Theoretical frameworks mobilised}

\caveat{The work declares no theoretical affiliation. What follows are
frameworks a critical study can apply productively, with an indication of what
in the work invites them --- not an attribution of intent to the author.}

\begin{description}[style=unboxed,leftmargin=0pt]
\item[Trauma and intergenerational transmission.] The creature figures pain
that crosses generations without being worked through. The
\textit{Historical Note} states it outright: ``\textit{the pain that crosses
generations, that no one processes and that is passed on like a contagion, is
the only true creature}''. This is the reading key the work offers of itself.

\item[Total institutions.] The Colony is described with the traits Goffman
identified in total institutions: collapse of the barriers between spheres of
life, mortification of the self, an economy of privileges. The ward of the
silenced intensifies that regime.

\item[Necropolitics and disposable lives.] The work insists that the Colony's
dead were ``poor, Black, epileptic, single mothers, homosexuals''. The
criterion for admission is social, not clinical --- which shifts the reading
from psychiatric horror to political horror.

\item[Archive, power and gap.] The chain of custody with documentary gaps
(DOC-20) and the archive incident reports (DOC-23) thematise the archive as an
instance of power and documentary silence as additional violence.

\item[Reception and the implicated reader.] The work constructs a model reader
and names them. The entire CONT family is a reception-theory device put to
work rather than discussed.
\end{description}

% =============================================================================
\section{Historical basis and ethical stance}
\label{sec:ethics}

The work's \textit{Historical Note} is methodologically careful, and this is
worth recording because it is rare in the genre:

\begin{itemize}
\item It states that Joaquim, the creature and the narrated institution are
fictional.
\item It distinguishes two historical referents the fiction might blur: the
1915 refugee corral (Alagadiço) and the Senador Pompeu one (1932--1933).
\item It states explicitly that ``\textit{the totals of inmates and deaths are
not presented as precise}''.
\item It attributes the documentary basis for the Barbacena Colony to a named
source (Arbex).
\item It declares the supernatural horror to be a metaphor, not a thesis.
\end{itemize}

The volume also carries a reader warning page, preceding the fiction, listing
sensitive content, distinguishing fiction from history and giving a crisis
line. It is paratext: it precedes the fictional frame and does not interrupt
immersion.

% =============================================================================
\section{Critical analysis}
\label{sec:critique}

\subsection{Where the work is strong}

\begin{enumerate}
\item \textbf{Form carries the argument.} Fragmentation, the asymmetry between
parts and the route network are not decorative: they formally reproduce the
impossibility of reconstruction and the duration of confinement.

\item \textbf{The reader's implication is executed, not announced.} The
numbered chain ends in a vacant position the reader occupies. It is a simple
device, rigorously built.

\item \textbf{Documentary honesty is explicit.} The work separates what it
invents from what it documents, and declines precision it does not have.

\item \textbf{The induction prose is technically competent.} The CONT-04
inversion is the best example: it removes the stimulus and keeps the effect.
\end{enumerate}

\subsection{Where the work is debatable}

\begin{enumerate}
\item \textbf{The supernatural metaphor may soften the historical horror.} By
giving the suffering a fantastic cause, the work risks displacing political
responsibility onto an entity. The \textit{Historical Note} anticipates the
objection and answers it; whether the answer suffices is an open critical
question, and the most productive one the work offers.

\item \textbf{Editor's notes interrupt the induction.} Thirteen fragments
carry footnotes declaring the device to be formal. Two land at the peak of the
effect --- CONT-03 closes on ``\textit{For your body, she is real}'' and the
note follows by stating that the contamination is textual. The information is
necessary; its placement is debatable. Consolidating the notes into a single
endnote would preserve both functions.

\item \textbf{The density of second person may saturate.} Twenty-seven
fragments sustain the ``you''. In continuous linear reading the procedure
tends towards the predictable --- a risk route-reading mitigates but the
printed edition does not remove.
\end{enumerate}

% =============================================================================
\section{Verified limitations}
\label{sec:limits}

\begin{enumerate}
\item \textbf{The three editions are not the same book.} Ten fragments diverge
structurally beyond what language difference explains. The extreme cases:
CONT-05 has 163 lines in PT and EN and 63 in ZH; CONT-03 has 174 in PT against
102 in EN and ZH; DOC-26 has 110 in PT and EN against 56 in ZH. In the
trilingual edition the three versions appear side by side, making the
divergence visible to the reader. \textbf{This is the most serious limitation
for any claim to being an international reference work.}

\item \textbf{Cover art is below print standard.} The cover and back-cover
files are 1,046 $\times$ 1,504 px, yielding 163 DPI on the printed panel; the
required standard is 300 DPI.

\item \textbf{The back cover carries a placeholder barcode.} The element
embedded in the artwork declares ISBN 978-65-01-23456-7, which is not the
work's ISBN.

\item \textbf{The corpus is not under version control.} The work's source
files have no commit history: there is no diff and no way to revert a bad edit.
\end{enumerate}

% =============================================================================
\section{Validation status}
\label{sec:validation}

This section exists to prevent the dossier being read as an endorsement.

\begin{center}\small
\begin{tabular}{@{}p{6.5cm}l@{}}
\toprule
\textbf{Dimension} & \textbf{Status} \\
\midrule
Build integrity of the 5 editions & verified \\
Route parity, source vs.\ printed & verified (648/648) \\
Convergence and absence of entrapment & verified \\
Fragment integrity (SHA-256) & verified \\
\midrule
Academic peer review & \textbf{not carried out} \\
External professional editorial review & \textbf{not carried out} \\
Historiographic verification by a specialist & \textbf{not carried out} \\
Native review of the EN and ZH translations & \textbf{not carried out} \\
Critical reception & \textbf{none} (work unpublished) \\
Impact & \textbf{not measurable} \\
\bottomrule
\end{tabular}
\end{center}

\caveat{The first four rows are \textbf{internal technical} verifications:
they say the object is well built, not that it is good. None substitutes for
critical judgement. The last six have not been carried out, and this dossier
does not anticipate them. Any claim of merit, originality or impact appearing
in material derived from this text should be treated as unfounded until
external validation exists.}

% =============================================================================
\section{Research agenda}

Questions the work opens that a study might take up:

\begin{enumerate}
\item Does the supernatural metaphor amplify or attenuate the testimonial
force of a narrative about documented atrocity?
\item What effects does the route network produce on the reading experience,
measured empirically rather than postulated?
\item Does second-person address sustained across 27 fragments produce
implication or reactance in readers?
\item How does divergence between the three editions alter the work as
received in each language --- and what does that say about the concept of a
``trilingual edition''?
\item How do the protective paratexts (warning page, editor's notes) interact
with the immersion device the work builds?
\end{enumerate}

% =============================================================================
\section*{Works referred to}

\begin{itemize}\small
\item ARBEX, Daniela. \textit{Holocausto Brasileiro}. Geração Editorial, 2013.
\item CALVINO, Italo. \textit{If on a winter's night a traveller}. 1979.
\item CORTÁZAR, Julio. \textit{Hopscotch}. 1963.
\item DANIELEWSKI, Mark Z. \textit{House of Leaves}. Pantheon, 2000.
\item GOFFMAN, Erving. \textit{Asylums}. 1961.
\item RAMOS, Graciliano. \textit{Barren Lives}. 1938.
\item STOKER, Bram. \textit{Dracula}. 1897.
\end{itemize}

\vfill
\noindent\rule{\linewidth}{0.4pt}\\
{\footnotesize Dossier produced as critical apparatus. On caveats regarding
the project's metrics and claims, see \texttt{CORRIGENDUM.md}.}

\end{document}
